\subsection{The Hahn-Banach theorem (geometric form)}

THEOREM 1 (Hahn-Banach). --- Let A be an open convex non empty set of the topological vector space E and let M be a non-empty linear variety which does not meet A. Then there exists a closed hyperplane H which contains M and does not meet A.

By translation the problem can be reduced to the case \(0 \in A\) , so that A is absorbent. Let p be the gauge of A (II, p.~20) so that A is the set of points \(x \in E\) such that \(p(x) < 1\) . On the other hand, let V be the vector subspace of E generated by M; thus M is a hyperplane in V that does not contain 0, and hence there is a unique linear form f on V such that M is the set of points \(y \in V\) for which \(f(y) = 1\) . The hypothesis \(M \cap A = \varnothing\) implies therefore that for all \(y \in V\) for which \(f(y) = 1\) , we have \(p(y) \geqslant 1\) ; as f and p are positively homogeneous we have \(f(y) \leqslant p(y)\) for all \(y \in V\) such that \(f(y) > 0\) ; finally as \(p(y) \geqslant 0\) for all \(y \in V\) , we see that \(f(y) \leqslant p(y)\) for all \(y \in V\) . By the analytical form of the Hahn-Banach theorem (II, p.~22, th. 1) there exists a linear form h on E which extends f and is such that, for all \(x \in E\) , \(h(x) \leqslant p(x)\) . Let H be the hyperplane in E with the equation \(h(x) = 1\) . Clearly \(H \cap V = M\) and \(H \cap A = \varnothing\) . On the other hand the complement of H in E contains the open non-empty set A, therefore H is closed in E (I, p.~11, corollary).

Q.E.D.

Remarks. --- 1) When \(0 \in \mathbf{M}\) , th. 1 can be stated as follows: there exists a continuous linear form in E, such that \(g(x) = 0\) in M and \(g(x) > 0\) in A (II, p.~8, prop. 4).

\begin{enumerate}
\def\labelenumi{\arabic{enumi})}
\setcounter{enumi}{1}
\tightlist
\item
  If we apply theorem 1 to the case where E carries the finest locally convex topology (II, p.~25, Example 2), and if, for the sake of simplicity, we suppose that \(0 \in A\) , then we get the following result (that superficially does not involve topology): if A is an absorbent convex set in the real vector space E and if M is a non-empty linear variety that does not meet A, then there exists a hyperplane H containing M and such that A lies on one side of H. This result is not valid for every convex set A (II, p.~65, exerc. 5).
\end{enumerate}

